\section{Определение квадратичной формы. Определение полярной билинейной формы}

\begin{definition}[Квадратичная форма и полярная билинейная форма]
Пусть $V$ --- линейное пространство над $\mathbb{R}$, $f(x, y)$ --- симметричная билинейная форма.

Функция $q(x) = f(x, x)$ называется \textbf{квадратичной формой}.

Исходная билинейная форма $f$ называется \textbf{полярной} для $q$. Связь между ними (формула поляризации):
\[
f(x, y) = \frac{1}{2}\bigl(f(x+y, x+y) - f(x, x) - f(y, y)\bigr).
\]

В координатах:
\[
f(x, x) = \sum_{i,j=1}^{n} a_{ij} \xi_i \xi_j = \sum_{i=1}^{n} a_{ii} \xi_i^2 + 2 \sum_{i < j} a_{ij} \xi_i \xi_j.
\]
\end{definition}
